\section{Introduction}
\label{sec:introduction}

eBPF is a technology for safely extending operating system kernels. It lets applications execute custom logic in the kernel for tasks such as packet processing~\cite{katran} and system observability~\cite{bpftrace}. Execution efficiency is an important requirement for these extensions. For example, on a network path that processes each packet, the eBPF program runs for every packet, so its execution cost affects CPU overhead and packet-processing capacity.

However, the existing eBPF compilation pipeline fails to select efficient implementations available on the target CPU for some computations. A userspace compiler first produces eBPF bytecode. After the kernel verifier~\cite{verifier} checks it, the in-kernel eBPF just-in-time (JIT) compiler generates native code. Some computations are decomposed into multiple basic operations in the bytecode. The current JIT mainly translates instructions individually and does not recombine these sequences into corresponding native operations. For example, a 64-bit rotation by a constant is represented as shifts and bitwise OR in the bytecode. The existing x86-64 JIT translates these operations separately, producing multiple native instructions. Direct native compilation can instead perform the same computation with a single rotate instruction.

Existing optimization approaches face different limitations in using these native implementations. Userspace optimizers such as K2~\cite{xu2021synthesizing} improve programs by rewriting existing eBPF instructions, but the resulting bytecode still goes through the existing JIT. Bytecode rewriting alone therefore cannot add native instruction-selection rules that the JIT lacks. Extending the JIT directly can add these rules, but requires maintaining and reviewing the corresponding pattern recognition, transformations, and code generation in the kernel. These transformations run after verification, so errors can also invalidate the verified safety conditions. Shahinfar et al.~\cite{shahinfar2026fun} propose native optimization in userspace after verification. This moves optimization code out of the kernel but requires trusting the userspace optimizer's translation.

To avoid adding a complex optimizer to the trusted code, we separate optimization selection from safety checking. The same local computation can have an eBPF implementation and a native implementation. The former is used for safety analysis of the complete program, and the latter for execution. Both must exhibit the same program-observable behavior for the safety analysis to apply to the code that executes.

We present \tool, an extensible framework for eBPF compilation that enables more efficient native implementations while reusing the existing verifier for safety analysis. A userspace pattern recognizer specifies operations and their parameters using extended instructions called \kinsns. The kernel checks their usage conditions and expands them into existing eBPF instructions, which the verifier analyzes with the rest of the program. After verification succeeds, the JIT generates code using semantically equivalent native implementations. We organize equivalence proofs around a shared operation specification and demonstrate a rotation proof in Lean~4~\cite{moura2021lean}. New recognition rules can reuse existing \kinsns. New operations reuse the framework's mechanisms for verification and code generation.

We implement \tool in Linux for x86-64 and ARM64, using kernel modules to provide seven categories of local optimizations involving computation, control flow, and memory access. With the stock Linux eBPF JIT as the baseline, the geometric mean microbenchmark speedups on the two architectures are 1.24$\times$ and 1.22$\times$, respectively. Throughput increases by 7.4\% for Cilium on x86-64 and by 7.3\% for Katran on ARM64. These results show that the framework's local optimizations improve the execution efficiency of computational microbenchmarks and the packet processing throughput of real applications.

\noindent\textbf{Contributions.}
We make the following contributions:
\begin{itemize}
    \item We investigate performance differences between eBPF and direct native compilation, and use code comparisons to identify local optimization opportunities that the existing pipeline leaves unused.
    \item We present \tool, connecting userspace pattern recognition to in-kernel native implementations via extended instructions, and methods to reuse existing safety analysis and prove semantic equivalence.
    \item We implement the framework and seven categories of local optimizations on x86-64 and ARM64, and evaluate their performance and overhead using microbenchmarks and packet-processing applications.
\end{itemize}